\documentclass[11pt,a4paper,oneside,openany]{book}
\usepackage{amsmath,amssymb,amsthm,mathrsfs}
\usepackage[margin=24mm,headheight=15pt]{geometry}
\usepackage[no-math]{fontspec}
\setmainfont{Noto Serif}
\setsansfont{Noto Sans}
\setmonofont[Scale=0.82]{DejaVu Sans Mono}
\usepackage{polyglossia}
\setdefaultlanguage{russian}
\setotherlanguage{english}
\usepackage{xcolor,enumitem,array,tabularx,longtable,booktabs}
\usepackage{tikz}
\usetikzlibrary{arrows.meta,positioning}
\usepackage{xurl}
\usepackage[unicode,colorlinks=true,linkcolor=blue!55!black,urlcolor=blue!55!black,citecolor=blue!55!black]{hyperref}
\usepackage{bookmark}
\usepackage{titlesec}
\titlespacing*{\chapter}{0pt}{12pt}{18pt}
\titlespacing*{\section}{0pt}{2ex plus .5ex minus .2ex}{.8ex plus .2ex}
\titlespacing*{\subsection}{0pt}{1.5ex plus .3ex minus .2ex}{.6ex plus .2ex}
\setlength{\emergencystretch}{3em}
\setlist{itemsep=3pt,topsep=5pt,leftmargin=*}
\setlist[description]{style=nextline,leftmargin=0pt,labelindent=0pt,labelwidth=0pt,labelsep=0pt}
\raggedbottom
\newcolumntype{Y}{>{\raggedright\arraybackslash}X}
\newcolumntype{P}[1]{>{\raggedright\arraybackslash}p{#1}}
\renewcommand{\arraystretch}{1.18}
\setcounter{tocdepth}{1}
\newcommand{\C}{\mathbb C}
\newcommand{\R}{\mathbb R}
\newcommand{\N}{\mathbb N}
\newcommand{\M}{\mathcal M}
\newcommand{\D}{\mathbb D}
\newcommand{\diam}{\operatorname{diam}}
\newcommand{\dist}{\operatorname{dist}}
\newcommand{\Comp}{\operatorname{Comp}}
\newcommand{\modu}{\operatorname{mod}}
\newcommand{\repo}[1]{\texttt{\path{#1}}}
\newcommand{\criterion}[1]{\par\smallskip\noindent\textbf{Критерий перехода.} #1\par\smallskip}
\newtheorem{exercise}{Задача}[chapter]
\newtheorem{proposition}{Утверждение}[chapter]
\theoremstyle{remark}
\newtheorem*{warningnote}{Замечание о статусе}
\title{Комплексная динамика и теория групп\\[6pt]
\large Программа работы над первой теоремой}
\author{}
\date{6 ноября 2026 года\\Нагрузка: 12--15 часов в неделю\\Версия 2.2}
\hypersetup{pdftitle={Комплексная динамика и теория групп: программа работы над первой теоремой},pdfauthor={}}
\begin{document}
\frontmatter
\maketitle
\chapter*{Как пользоваться программой}
\addcontentsline{toc}{chapter}{Как пользоваться программой}
Цель первого цикла --- получить полное доказательство одного содержательного
утверждения, установить его отношение к известным результатам и подготовить
исследовательскую записку для внешней проверки.

В первые две недели выбирается точная постановка. Следующие десять недель
отводятся на доказательство, примеры и текст. Двенадцать недель --- срок
исследовательской попытки; существование нового результата за этот срок
заранее неизвестно. Если постановка оказывается известной или неверной,
решение о её изменении принимается на ближайшем контрольном этапе.

Начать с глав~1--3 и файла \repo{first_12_weeks.md}.
Основная тема --- бисеты. Перед основным циклом доказательства
выполняется проверка предполагаемого научного вклада по
\repo{templates/novelty_card.md}. Образец первых действий ---
\repo{first_days.pdf}; завершённая проверка B1-R ---
\repo{novelty_audit.pdf}. Формулировки и аргументы ведутся в \repo{research/}.
Литература выбирается по зависимостям текущего доказательства.
Каждая неделя должна оставлять письменный математический результат.
\begingroup\footnotesize
\tableofcontents
\endgroup
\mainmatter
\chapter{Задача первого исследовательского цикла}

\section{Результат, на который направлена работа}
Подготовить записку с одной основной теоремой, полным доказательством,
проверенным примером и точным сравнением с ближайшей литературой.
Основная теорема должна решать явно сформулированную трудность:
давать новую конструкцию, оценку, класс применимости или препятствие
для существующего метода.

Учебные леммы и воспроизведения известных доказательств включаются в
рабочие записи с указанием источника. Статус основного результата
пересматривается после поиска предшественников и внешнего обсуждения.

Первый цикл занимает 12 учебных недель, то есть 144--180 часов.
Контрольные этапы --- конец недель~2, 4, 8 и~12.
Выбор задачи начинается в первый рабочий день.

\section{Что уже есть в репозитории}
\begin{description}
\item[Достаточный топологический критерий.]
\repo{Mlc/ParameterComponentApproximation.lean}:
буферное условие для убывающих компактных приближений влечёт локальную
связность предела. Для множества Мандельброта соответствующий вход
передаётся явно.
\item[Форма конечного покрытия.]
\repo{Mlc/FiniteComponentCriterion.lean}:
поле \repo{LocalPiece.control} всё ещё содержит
\(\forall N\,\exists L\) и включение в компоненту.
Для него требуется отдельное математическое обоснование.
\item[Доступные конструкции.]
Внешние орбитальные приближения, рациональные прямоугольники, интервальные
оценки и ловушки. Начальные модули ---
\repo{Mlc/CertifiedOrbitApproximation.lean},
\repo{Mlc/CertifiedTrappingRegions.lean},
\repo{Mlc/FiniteTrappingRegions.lean}.
\item[Установленные ограничения моделей.]
\repo{Mlc/GreenSublevelIntersectionCounterexample.lean} содержит
контрпример к прежнему универсальному утверждению.
\repo{Mlc/Literature/ModelAudit.lean} выявляет последствия подмены
конформного модуля гауссовой площадью.
\end{description}
Для каждого используемого результата записать точную декларацию, автора,
дату и источник по истории проекта. Условные входы раскрыть до определений.

\section{Проверки постановки}
До длинного доказательства выполнить четыре действия:
\begin{enumerate}
\item Предъявить объект, удовлетворяющий всем гипотезам, и вычислить его
      входные данные.
\item Проверить, что гипотезы можно установить без предполагаемого вывода.
\item Найти ближайшую известную теорему и проверить возможность прямого
      применения к выбранному классу.
\item Назвать один переход, для которого в проекте действительно
      отсутствует доказательство.
\end{enumerate}
Наличие слова «конечный» в названии структуры проверяется по типам её полей.
Для алгоритма отдельно проверяются представление входа, завершаемость
и смысл выдаваемого ответа.

\section{Подготовка по возникающей зависимости}
Для недостающей леммы записать: место её использования, точную формулировку,
источник и требуемый объём разбора. Сначала попытаться доказать нужный
частный случай. Если требуется крупная внешняя теорема, привести её
формулировку и проверить гипотезы.

Механический ознакомительный проход --- медленное чтение и воспроизведение
обозначений и полных формулировок. Он помогает освоить форму незнакомого
фрагмента. В этом цикле такой проход занимает до 30 минут внутри
читательского бюджета и заканчивается указанием нужной леммы.
Непрерывное чтение Любича после с.~133 временно снимается с расписания.
Михайлов--Пасси, Хэтчер и Милнор подключаются при появлении соответствующих
зависимостей.

\chapter{Будущая работа и выбор узкой задачи}

\section{Рабочая тема}
Предлагаемое направление по умолчанию:
\begin{quote}
\textbf{Геометрические критерии локальной связности и проверяемые конечные
аппроксимации в комплексной динамике.}
\end{quote}
После пилота зафиксировать класс объектов, формат сертификата и основную
теорему. Обсудить постановку с руководителем.

\section{Три варианта специализации}
\begin{description}
\item[A. Геометрия и сертифицированные приближения --- основной кандидат.]
Построить конечные сертификаты включения в компоненту для заданного класса
геометрических конфигураций. Получить оценки размеров, глубины и точности
через геометрические запасы.
\item[B. Алгебраическая динамика конечного посткритического типа.]
Выбрать одно утверждение Bartholdi--Dudko о заданном классе разветвлённых
накрытий: алгоритм с доказательством корректности или структурную теорему.
\item[C. Аналитическая ренормализация.]
Восстановить один геометрический переход из principal nest или a priori bounds,
затем сформулировать и проверить усиление.
\end{description}
После 36-й недели оставляется одна основная ветка. Вторая может занимать
не более одной короткой сессии в две недели; третья откладывается.

\section{Что именно исследовать в ветке A}
Текущая схема использует
\[
O_N=\{c:|c|\le2,\ |f_c^n(0)|\le2\ (n\le N)\},\qquad
\bigcap_N O_N=\M.
\]
Для \(x\in\M\) обозначим
\[
C_N(x,r)=\Comp_x\bigl(O_N\cap\overline B(x,r)\bigr).
\]
Нужное достаточное условие имеет порядок кванторов
\begin{equation}\label{eq:buffer}
\forall x\in\M\ \forall\varepsilon>0\ \exists\,0<\delta<r<\varepsilon\quad
\forall N\ \exists L\ge N:
O_L\cap\overline B(x,\delta)\subseteq C_N(x,r).
\end{equation}
Радиусы выбираются \emph{до} \(N\), а \(L\) может от него зависеть.
Достаточность этого условия доказана в
\repo{Mlc/ParameterComponentApproximation.lean}.

В \repo{Mlc/FiniteComponentCriterion.lean} поле
\repo{LocalPiece.control} содержит \(\forall N\,\exists L\) и включение
в компоненту. Задача пилота --- задать конечные геометрические данные,
из которых следует такое включение в выбранном классе.

\textbf{Спецификация пилота.}
\begin{quote}
Для явно определённого класса геометрических конфигураций \(\mathcal C\)
построить формат рационального сертификата \(W\), проверяющий включение
\(O_L\cap\overline B(x,\delta)\subseteq C_N(x,r)\), и доказать:
(i) корректность проверки; (ii) существование сертификата при указанных
геометрических запасах; (iii) явную зависимость требуемой глубины и точности
от этих запасов.
\end{quote}
Начать с фиксированных \(x,N,r,\delta,L\). Переход к общим
кванторам~\eqref{eq:buffer} выделить в отдельную теорему.

Если использовать граф клеток, нужно доказать связность самих клеток,
достоверность их пересечений и нахождение в нужном внешнем множестве.
Для полуалгебраической постановки конечного уровня рассмотреть точные
методы реальной алгебраической геометрии и оценить их сложность.

\section{Проверка силы достаточного условия}
Локальная связность предела не влечёт буферное свойство произвольных
компактных внешних приближений. Контрпример:
\[
S=\{0\}\subset\R,\qquad
O_N=\{0\}\cup\{1/k:k\ge N+1\}.
\]
Все \(O_N\) компактны, убывают и имеют пересечение \(S\), которое локально
связно. Однако компонента нуля в \(O_N\cap[-r,r]\) равна \(\{0\}\):
подмножества прямой, содержащие два различных числа и связные, должны
содержать интервал, а этот набор счётен. Для любых \(L\) и \(\delta>0\)
в \(O_L\cap[-\delta,\delta]\) есть положительное \(1/k\). Значит, нужное
буферное включение не выполняется.

\section{Проверка новизны}
На 24-й и 36-й неделях для двух кандидатов составить таблицу: точная
формулировка; ближайшая теорема; что изменено; почему это не непосредственное
следствие; кто может проверить; минимальный контрпример к чрезмерному усилению.
Проверить доказательства и библиографии ближайших работ, затем передать
сравнение внешнему рецензенту.

\section{Каркас будущей рукописи}
\begin{enumerate}
\item Задача, точная общность и отличие от известных результатов.
\item Необходимые аналитические и топологические основания с доказательствами.
\item Ограничения прежних моделей; авторство и история каждого результата.
\item Основная новая конструкция или теорема.
\item Примеры, границы применимости и независимые проверки.
\item Lean-реализация: соответствие математики определениям, доверенная база,
воспроизводимость; оставшиеся обязательства вынесены явно.
\end{enumerate}
Обзор из \repo{docs/} использовать как справочник при подготовке рукописи.

\chapter{Рабочая неделя и доказательство}

\section{Бюджет 12--15 часов}
Базовая неделя занимает 12 часов.
\begin{center}
\begin{tabularx}{\textwidth}{@{}P{0.38\textwidth}P{0.1\textwidth}Y@{}}
\toprule
Работа & Время & Результат\\
\midrule
Основная лемма: три сессии & 6 ч & Письменные попытки, доказанный переход
или контрпример\\
Литература по текущему препятствию & 2 ч & Точная внешняя теорема,
проверка гипотез и сравнение\\
Чистовой математический текст & 2 ч & Исправленная формулировка
и связный фрагмент доказательства\\
Примеры, вычисления или Lean & 1 ч & Воспроизводимая проверка
одного спорного места\\
Обсуждение и журнал & 1 ч & Одно главное препятствие и задача
следующей недели\\
\bottomrule
\end{tabularx}
\end{center}
При 15 часах добавить два часа к доказательству и час к проверкам
или внешнему обсуждению. В первые две недели действует бюджет отбора:
4 часа --- близкие теоремы и доказательства, 3 часа --- попытки сведения,
3 часа --- формулировка и сравнение, 1 час --- примеры,
1 час --- обсуждение и журнал. После обоснования дополнительного вопроса
применяется основная таблица. Семинар и переписка
с исследователем входят в общий бюджет.

Параллельных обязательных курсов в этом цикле нет. Чтение Любича,
Михайлова--Пасси, Хэтчера и Милнора оплачивается из двух часов литературы.
Увеличение этого бюджета на конкретной неделе записывается вместе
с названием закрываемой зависимости.

\section{Одна исследовательская сессия}
\begin{enumerate}
\item Назвать один переход, который должен быть получен за сессию.
\item Выписать все доступные гипотезы и проверить простейший пример.
\item Работать над аргументом 60--90 минут. При остановке записать
      последнее корректное утверждение и точный пробел.
\item Последние 15--20 минут посвятить тексту: доказанный шаг, неудачная
      оценка, контрпример или уточнение постановки.
\end{enumerate}
Если найдено неверное утверждение, сохранить его вместе с контрпримером.
Рабочие файлы первой темы перечислены в \repo{research/README.md}.

\section{Разбор нужного известного доказательства}
Сначала выделить используемую лемму и попытаться доказать нужный частный
случай. Затем прочитать источник, проверить условия и закрыть книгу.
Восстановить аргумент письменно; через неделю повторить ключевой переход
и проверить изменённый пример.

Для внешнего результата сохраняются автор, версия, номер теоремы,
страница и проверка гипотез. Полное воспроизведение крупной теории
назначается отдельной задачей только при необходимости для основного
доказательства.

\section{Состояния работы}
В журнале различать: \emph{кандидат}, \emph{в работе},
\emph{доказано на бумаге}, \emph{проверено извне},
\emph{опровергнуто}, \emph{сведено к известному}.
Новизна фиксируется отдельно: \emph{не проверена},
\emph{найден предшественник}, \emph{элементарное следствие},
\emph{добавление предварительно обосновано},
\emph{сравнение обсуждено со специалистом}.
Научный вес описывается отдельным абзацем с примером предполагаемой пользы.
Полученный отзыв записывается дословно по существу;
отсутствие ответа остаётся отсутствием проверки.

\section{ИИ и Lean}
Использовать ИИ для одной конкретной леммы, поиска ошибки, проверки
кванторов или подготовки формализации. Предложенный аргумент получает
статус после проверки каждого перехода и используемых источников.
Автора идеи, текст доказательства и способ проверки отмечать в карточке.

Формализовать устойчивое бумажное утверждение. В базовом бюджете
на Lean и вычисления приходится один час в неделю; при 15 часах
можно выделить до двух. После формализации сверить определения,
проверить зависимости через \repo{#print axioms} и собрать модуль.
Техническая задержка переносится в отдельный список.

\section{Правило остановки}
После двух сессий без продвижения выписать препятствие как отдельную
лемму. После четырёх --- проверить модельный случай, источник или получить
предметный разбор. После двух недель с тем же пробелом уменьшить общность
или изменить конструкцию. Решение записать с математической причиной.

\chapter{Первые двенадцать учебных недель}

\section{Две недели диагностики и восстановления}
В первые две недели выбрать четыре из восьми задач ниже, обязательно взяв
одну топологическую и одну аналитическую. На каждую выделить до 90 минут
первой самостоятельной попытки, затем до двух часов разбора и исправления.
Остальные задачи остаются ориентирами на квартал. Весь список не требуется
решить за две недели.

\begin{exercise}[Локальная связность и гипотезы]
Дать определение локальной связности подпространства и критерий через
связные окрестности внутри произвольной окрестности. Привести два связных
компактных множества в плоскости с несвязным непустым пересечением.
Объяснить, почему принадлежность центральной точки пересечению не помогает.
\end{exercise}

\begin{exercise}[Критическая орбита]
Доказать, что ограниченность критической орбиты \(f_c(z)=z^2+c\)
влечёт \(|c|\le2\) и \(|f_c^n(0)|\le2\) для всех \(n\).
Отсюда вывести компактность \(\M\) через непрерывность конечных орбитальных
полиномов. Если начинается геометрический рост, указать его постоянный
множитель, больший единицы.
\end{exercise}

\begin{exercise}[Настоящий модуль]
Для \(0<r<R\) вычислить экстремальную длину семейства кривых,
соединяющих две граничные окружности кольца \(A(r,R)\).
Получить \(\modu A=(2\pi)^{-1}\log(R/r)\).
Сначала объявить нормировку \(\sup_\rho L_\rho(\Gamma)^2/A_\rho\),
затем доказать верхнюю и нижнюю оценки. Проверить масштабирование.
\end{exercise}

\begin{exercise}[Компактные вложенные множества]
Пусть непустые компакты \(K_n\) в метрическом пространстве убывают и
\(\bigcap K_n=\{x\}\). Доказать \(\diam K_n\to0\).
Отдельно разобрать, где используется компактность \(K_0\), и построить
контрпример без неё. Версия с окрестностями: для любой открытой \(U\ni x\)
найдётся \(n\), для которого \(K_n\subset U\).
\end{exercise}

\begin{exercise}[Лемма Шварца--Пика]
Нормировать автоморфизмами диска произвольное голоморфное
\(h:\D\to\D\) и вывести
\[
\frac{|h'(z)|}{1-|h(z)|^2}\le\frac1{1-|z|^2}.
\]
Разобрать случай равенства. Указать, что неувеличение гиперболического
расстояния ещё не даёт равномерного сжатия с постоянной меньше единицы.
\end{exercise}

\begin{exercise}[От колец к сжатию]
Сформулировать последовательное неравенство Грётша с точными условиями
взаимного расположения колец. Вывести один критерий сжатия вложенных
топологических дисков, явно указав нужную теорему о конформном модуле.
Проверить аргумент на концентрических дисках и сравнить с гауссовой площадью.
На диагностике достаточно локализовать недостающую аналитическую лемму;
её полное освоение запланировано отдельно.
\end{exercise}

\begin{exercise}[Кванторы условного корня]
Развернуть \repo{UniformOuterBuffer} на бумаге и воспроизвести из него
связную окрестность точки предельного компакта. Не использовать корневую
теорему как посылку. Отметить, что изменится, если \(\delta\) зависит от \(N\).
\end{exercise}

\begin{exercise}[Алгебраическая ветка]
Дать определения нижнего центрального ряда \(\gamma_n(G)\) и размерного
ряда относительно явно выбранного коэффициентного кольца. Для свободной
абелевой группы вычислить оба ряда. Изобразить, как подъём петли при
\(z\mapsto z^2:\C^*\to\C^*\) действует на фундаментальной группе.
Объяснить, какие дополнительные данные нужны для перехода к бисету.
\end{exercise}

\criterion{Если выбранные базовые задачи восстанавливаются самостоятельно,
не повторять весь вводный курс: переходить к новым разделам и использовать
освободившиеся сессии для исследовательских примеров. Если возникли пробелы,
вставить двухнедельный ремонт конкретной зависимости, не начинать книгу сначала.}

\section{Расписание оставшихся недель}
Подробные действия и файлы результатов перечислены в
\repo{first_12_weeks.md}. Здесь указана математическая логика.

\begin{table}[!tp]
\centering\small
\caption{Первые двенадцать учебных недель}
\begin{tabularx}{\textwidth}{@{}P{0.07\textwidth}P{0.37\textwidth}Y@{}}
\toprule
Нед. & Материал & Один основной результат\\
\midrule
1 & Диагностика: орбиты и топология & Две попытки без подсказок; одна
исправленная карточка с доказательством\\
2 & Аудит определений и две оставшиеся диагностические задачи & Сопоставление
модуля кольца и weighted area; карта условного корня\\
3 & Любич, §6, с.~119--132: выбранные пункты & Полное вычисление модуля
круглого кольца и проверка нормировок\\
4 & Последовательное неравенство и вложенные компакты & Одно доказательство
сжатия в явно указанной общности; квартальный журнал ещё не нужен, только
первый четырёхнедельный отчёт\\
5 & Любич, §7, с.~133--145: начало & Вывод метрики диска выбранной
нормировки и инвариантности относительно автоморфизмов\\
6 & Продолжение §7 & Доказательство Шварца--Пика, случай равенства,
пример отсутствия равномерной строгой константы\\
7 & Компактность и компоненты из текущего проекта & Восстановление одного
перехода в доказательстве достаточности буфера\\
8 & Любич, §8, с.~146--151 & Точная формулировка граничного продолжения
Риманова отображения; пример, где продолжение не инъективно\\
9 & Любич, приложение §10, с.~163--172; выбранные леммы & Принцип минимума
для гармонической функции в нужной ограниченной области; условия на границе\\
10 & Вводный Бёттхер: §23, с.~312--329, один фрагмент & Функциональное
уравнение Green-функции из определения и нулевой уровень на ограниченной орбите\\
11 & Любич, §9, с.~152--162; пример замороженной модели & Сравнение полного
Green-подуровня и граф-разрезанного puzzle; конечный пример \(c=0\) или \(-1\)\\
12 & Повторение и письменная сборка & Связная записка с 3--5 проверенными
утверждениями, устный доклад и пересмотр паспорта\\
\bottomrule
\end{tabularx}
\end{table}

Страницы относятся к местной книге 2023 года. Порядок намеренно не
совпадает с последовательным чтением всех страниц: пропускаемые разделы
возвращаются тогда, когда понадобятся их результаты. В неделе~10 не
требуется за две сессии доказать весь Бёттхер и глобальную униформизацию.

\section{Первые законченные тексты}
К концу квартала ожидаются примерно 6--8 небольших проверенных доказательств,
3--5 разобранных контрпримеров и одна объединяющая записка ориентировочно на
8--15 страниц. Это ориентиры для масштаба, а не норматив объёма. Если одна
лемма потребовала три недели и полностью понята, она ценнее десятка
непроверенных «готовых» доказательств.

В записке удобно связать три темы: настоящий модуль кольца; компактное
сжатие и компоненты; граница перехода от динамики к параметрам. Крупные
внешние теоремы перечисляются явно. Записка проверяет владение материалом,
а не заявляет новый результат.

\section{Первое внешнее обсуждение}
В недели~4--8 подготовить одну страницу о своей подготовке и двух возможных
темах, приложить 2--4 страницы одной леммы. Самостоятельно обратиться к
исследователю или семинару с конкретным вопросом о посильной задаче.
Цель --- получить содержательную критику проверяемого аргумента.

\chapter{Маршрут на два года и развилки}

\section{Календарная модель}
В году 44 учебные недели и восемь резервных. Каждая четвёртая неделя включает
отчёт. Календарные даты заносятся в журнал при начале занятий.
При невыполненном критерии следующий этап переносится.

\section{Первый год: основания и выбор проекта}
\subsection*{Недели 1--12: проверяемая основа}
Материал и задания заданы в предыдущей главе. Особое внимание уделяется
нормировкам, компонентам, компактности и различию двух плоскостей.
\criterion{Вы можете доказать ключевые леммы из квартальной записки у доски
и объяснить, почему две отвергнутые модели репозитория не годятся.}

\subsection*{Недели 13--20: квазиконформная геометрия}
Любич, §§11--18, с.~175--229. Выписать определения и цепочку зависимостей;
полностью разобрать 2--3 доказательства.
Минимум: аналитическое и геометрическое определения qc-отображения,
искажение модуля, нормировка решения уравнения Бельтрами, смысл голоморфного
движения и его продолжения. MRMT и теорему продолжения можно сначала
использовать как явно названные внешние теоремы. Полное доказательство MRMT
выделить в отдельный блок, если оно нужно выбранной теме.

Вычисления: проверить радиальное растяжение \(z\mapsto z|z|^{\alpha-1}\)
при \(\alpha>0\), его дилатацию и действие на модуль кольца.
\criterion{Письменный расчёт qc-примера, точная теорема об искажении модуля
и одна схема склейки с обозначенными входными условиями.}

\subsection*{Недели 21--28: динамические основания и первый отбор тем}
Любич, §§19--24, с.~234--344, выборочно. Нужны классификация периодических
точек, нормальные семейства, заполненное множество Жюли, координата Бёттхера
у бесконечности и роль посадки рациональных лучей. На этом этапе различать
локальное построение координаты и её глобальное продолжение, а также
утверждения для отдельных лучей и для всей границы.

К концу 24-й недели составить два исследовательских паспорта и сопоставить
их с существующими декларациями проекта.
\criterion{Один полный аргумент из динамики и два точных кандидата на
исследовательский вопрос с ближайшими источниками.}

\subsection*{Недели 29--36: puzzle, polynomial-like и параметрический переход}
Обзор определений из §§25, 28, 31, 33--34, 38, 40, 42, 47 книги Любича.
Для полного разбора выбрать один механизм: компактное
сжатие корректных параметрических кусочков, конкретный конечный puzzle
или локальную polynomial-like конструкцию.

Нарисовать отдельно фазовый и параметрический графы. Подписать, что именно
движется с параметром, где использовано голоморфное движение и откуда
берётся связность параметрической части. При невозможности доказать шаг
указать точную литературу и гипотезы. Сравнить с \repo{ParameterPuzzleBasis}
из текущих \repo{Mlc/Literature/} модулей.
\criterion{Доклад 30--40 минут, один проверенный механизм и выбор единственной
основной ветки A, B или C после внешнего обсуждения.}

\subsection*{Недели 37--44: восьминедельный исследовательский пилот}
У пилота одна центральная лемма и один класс примеров. Каждые две недели
должен появляться проверяемый объект: доказательство, опровергнутая гипотеза,
корректный сертификат или уменьшенная постановка. К окончанию пилота написать
результат как короткую статью: постановка, связь с известным, определения,
основное утверждение, доказательство, пример и ограничения.
\criterion{Законченная записка по пилоту и задача на следующий цикл.}

\section{Второй год: один вклад и законченная рукопись}
\begin{table}[!htp]
\centering\small
\caption{Второй год: этапы и результаты}
\begin{tabularx}{\textwidth}{@{}P{0.13\textwidth}P{0.4\textwidth}Y@{}}
\toprule
Недели & Работа & Условие перехода\\
\midrule
45--56 & Уточнение основной конструкции; проверка новизны; исправление
семантики и один нетривиальный пример & Зафиксированы утверждение, точный
класс объектов, ближайший результат и согласованный внешний отзыв\\
57--68 & Главная лемма; оценки и контрпримеры к лишним обобщениям &
Есть полный письменный аргумент либо доказана причина отказа от него;
следующий цикл имеет изменённую постановку\\
69--76 & Независимая проверка; выбранная Lean-реализация; воспроизводимые
сертификаты & Замечания разобраны, проверенные части отделены от условных\\
77--84 & Сборка рукописи; введение переписывается под фактический результат &
В тексте нет неявных переходов между плоскостями, неподтверждённой новизны
и пустых гипотез под именем содержательной теоремы\\
85--88 & Два пробных доклада; ответы на возражения; решение о подаче или
третьем годе & Внешний математик понимает вклад, вы можете защитить ключевое
доказательство без ИИ; дальнейшие требования согласованы отдельно\\
\bottomrule
\end{tabularx}
\end{table}

На неделе~68 по результатам рецензии выбрать продолжение: исходная задача,
более узкая постановка или дополнительный исследовательский год.

\section{Алгебраическая развилка B}
Если эта ветка оказывается ближе вашим интересам, после общего квартала
заменить часть аналитического маршрута на 12--16 недель следующей цепочки:
\begin{enumerate}
\item Накрытия и подъём путей; фундаментальная группа проколотой сферы;
      действие монодромии. Самостоятельно вычислить случай \(z\mapsto z^2\).
\item Коммутирующие левое и правое действия; тензорное произведение бисетов;
      зависимость от базовой точки и выбранных путей.
\item Bartholdi--Dudko I~\cite{BDI}: определения и один ограниченный случай
      теоремы ван Кампена для бисетов.
\item Один пример конечного посткритического типа, например \(z^2-1\),
      с полностью описанными маркировками и действиями.
\item Затем только нужный фрагмент IV/V~\cite{BDIV}, если исследовательский
      вопрос действительно касается расширяющей структуры.
\end{enumerate}
Групповые ряды и гомологические инструменты изучать по зависимостям
выбранного утверждения.

\section{Развилка C: современная аналитическая геометрия}
После освоения основных инструментов выбрать один маршрут:
\begin{enumerate}
\item конечная ренормализуемость и Yoccoz;
\item quasi-additivity и одна лемма переноса модуля Kahn--Lyubich;
\item одна указанная гипотеза теоремы о спутниковых или Feigenbaum-параметрах.
\end{enumerate}
Сначала восстановить ограниченную известную цепочку, затем искать усиление.
По обзору Дудко~\cite{DudkoSurvey} составить карту зависимостей.
Работы о спутниковых параметрах и Feigenbaum-точках
дают отдельные классы локальной связности~\cite{Satellite,Feigenbaum}.
Работа Kahn--Kapiamba--Lyubich 2026 года относится к параболически ограниченной
примитивной ренормализации~\cite{KKL}.

\chapter{Проверка результата}

\section{Формулировка и аргумент}
\begin{enumerate}
\item Все объекты определены; порядок кванторов зафиксирован.
\item Гипотезы проверены на примере и устанавливаются без предполагаемого
      вывода.
\item У конечного входа указано представление чисел и структур;
      у оценки --- параметры и модель вычислений.
\item Для каждой леммы выписаны зависимости и проверены условия внешних теорем.
\item Различены свободная группа и группа точного действия,
      стабилизатор уровня и стабилизатор одной вершины.
\item Согласованы сторона действия, порядок произведения секций,
      переход к фактору и независимость от базиса.
\item Ранг вычислен точной арифметикой; конечные значения
      и утверждение для всех глубин имеют отдельные доказательства.
\end{enumerate}
Перечитать аргумент от заключения к предпосылкам. Затем изложить ключевую
лемму без рабочих записей и проверить удаление одной гипотезы.
К вычислениям приложить точные входные данные.

\section{Новизна и внешняя проверка}
Для ближайшей теоремы записать номер, страницу, гипотезы, вывод и отношение
к кандидату. Проверить нужный фрагмент доказательства, предшественников
и продолжения. На неделе~2 подготовить вопрос для обсуждения;
на неделях~4 и~8 --- конкретную лемму с доказательством.
Для замечания записать ответ и изменение текста.

Первая теорема закончена при наличии полного доказательства,
проверенного примера, сравнения с литературой и связной записки.
Внешний отзыв фиксируется отдельно с датой и объёмом проверки.
До его получения результат имеет статус авторской проверки.
Решение о публикации принимается после оценки вклада и разбора замечаний.

\chapter{Источники и порядок чтения}

\section{Рабочий набор литературы}

\begin{center}\small
\begin{tabularx}{\textwidth}{@{}P{0.22\textwidth}P{0.2\textwidth}Y@{}}
\toprule
Источник & Когда & Материал и результат\\
\midrule
Любич~\cite{LyubichBook} & 2 ч в неделю & Выбранные разделы конформной
геометрии; определение, формулировка и один пример\\
Михайлов--Пасси~\cite{MP} & 2 ч в неделю & Коммутаторы, групповые кольца,
размерные подгруппы; доказательство на 1--2 недели\\
Хэтчер~\cite{Hatcher} & Недели 1--8 & §§1.1, 1.3; §1.2 для вычисления
фундаментальной группы проколотой плоскости\\
Милнор~\cite{MilnorThird,MilnorNotes} & Недели 9--12 и 21--28 & Fatou--Julia,
локальная динамика, заполненное множество Жюли и внешние лучи\\
Некрашевич~\cite{NekraIMG,NekraBook} & После топологического блока &
Сначала §§2--3 обзора; монография по выбранной задаче\\
\bottomrule
\end{tabularx}
\end{center}

\section{Выборочный маршрут по Любичу}
Местная версия --- \repo{refs/Lyubich_book_vol_I-II_2023.pdf}.
Пагинация сверена по оглавлению.
\begin{enumerate}
\item §7, с.~133--145: текущий ознакомительный проход и выбранные доказательства.
\item §8, с.~146--151, и §10, с.~163--172: граничное продолжение и
      необходимые леммы теории потенциала.
\item §6, с.~119--132: возвращаться для доказательств о модуле.
      §9, с.~152--162: читать при работе с puzzle после топологического блока.
\item §§11--17: определения qc-геометрии, искажение модуля, MRMT,
      голоморфные движения; 2--3 выбранных доказательства.
\item §§28, 31, 33--34, 38, 40, 42, 47: материал исследовательской ветки
      после выбора конкретного механизма.
\end{enumerate}
Вводную динамику изучать по Милнору. К соответствующим пунктам §§19--24
Любича обращаться за деталями доказательства. §18 и остальные разделы
добавлять по зависимостям задачи.

\section{Доказательства из Михайлова--Пасси}
Первые ориентиры --- §1.1, §1.2, начало главы~2 и §4.1.
Определения размерных подгрупп брать над \(\mathbb Z\):
\[
\Delta=\ker\bigl(\epsilon:\mathbb Z[G]\to\mathbb Z\bigr),\qquad
D_n(G)=\{g\in G:g-1\in\Delta^n\}.
\]
На квартал выбрать 4--6 доказательств из следующей цепочки:
\begin{enumerate}
\item Вывести коммутаторные тождества и доказать характеристичность
      \(\gamma_n(G)\). Проверить включение
\[
\gamma_n(G)/\gamma_{n+1}(G)\subseteq
\operatorname{Z}\bigl(G/\gamma_{n+1}(G)\bigr).
\]
\item Изоморфизм \(G_{\mathrm{ab}}\simeq\Delta/\Delta^2\), заданный
      \(g\mapsto(g-1)+\Delta^2\); следствие \(D_2(G)=\gamma_2(G)\).
\item Нормальность \(D_n(G)\) и включение
      \([D_m(G),D_n(G)]\subseteq D_{m+n}(G)\). Использовать тождество
\[
[x,y]-1=x^{-1}y^{-1}\bigl((x-1)(y-1)-(y-1)(x-1)\bigr),
\]
      где \([x,y]=x^{-1}y^{-1}xy\).
\item Включение \(\gamma_n(G)\subseteq D_n(G)\) по индукции.
\item Одна лемма о свободной группе или следствие коммутаторного исчисления
      из §§1.1--1.2; выписать все используемые тождества.
\item Разложение Магнуса \(x_i\mapsto1+X_i\): вычислить коммутатор
      до степени~2 в некоммутирующих переменных.
\end{enumerate}
На каждое доказательство оформить карточку и проверить конкретную группу.
В главах о гомологиях и симплициальных методах выбирать результаты по
исследовательскому паспорту.

\section{Топологический и динамический блоки}
По Хэтчеру разобрать подъём путей и гомотопий и действие фундаментальной
группы на слое накрытия. Вычислить накрытие \(z\mapsto z^2\) на \(\mathbb C^*\)
и один пример для проколотой плоскости. Затем открыть определения
итерационной монодромии в обзоре Некрашевича.

Для Милнора использовать третье издание либо местный ранний конспект
\repo{refs/9201272v1.pdf}. В раннем конспекте выбрать §§3, 6, 15--18;
в третьем издании ориентироваться на названия разделов.
Первый результат --- одно доказательство о локальной динамике и один
разобранный переход к заполненному множеству Жюли или внешним лучам.

\section{Современные статьи: по одной цепочке}
Для выбранного результата выписать: класс параметров, ограничения
комбинаторики, геометрическую оценку, используемый параметрический переход
и заключение. Разбирать одну зависимость за раз. Карта источников ---
\repo{refs/README.md} и \repo{refs/mlc_sources.json}.

\section{Обновление литературы}
Раз в месяц выделять до 30 минут на страницы авторов и новые версии только
выбранных работ. В журнал заносить изменения утверждений и их влияние
на текущую задачу. Раз в квартал обновлять карту литературы.

Интернет-ссылки проверены 6 ноября 2026 года. В тексте фиксировать
использованную версию статьи.

\begin{thebibliography}{99}\small
\bibitem{LyubichBook}
M.~Lyubich. \emph{Conformal Geometry and Dynamics of Quadratic Polynomials,
vol.~I--II}. Online manuscript; локальная версия 2023 года.
\href{https://www.math.stonybrook.edu/~mlyubich/book.pdf}{Книга};
\href{https://www.math.stonybrook.edu/~mlyubich/papers/index.html}{страница автора}.

\bibitem{MP}
R.~Mikhailov, I.~B.~S.~Passi.
\emph{Lower Central and Dimension Series of Groups}.
Lecture Notes in Mathematics 1952, Springer, 2009.
\href{https://link.springer.com/book/10.1007/978-3-540-85818-8}{Издательская страница и оглавление}.

\bibitem{DH}
A.~Douady, J.~H.~Hubbard.
\emph{Exploring the Mandelbrot Set: The Orsay Notes}.
\href{https://pi.math.cornell.edu/~hubbard/OrsayEnglish.pdf}{Авторская публикация}.

\bibitem{DudkoSurvey}
D.~Dudko. \emph{On the MLC Conjecture and the Renormalization Theory
in Complex Dynamics}. arXiv:2512.24171v1, 2025.
\href{https://arxiv.org/abs/2512.24171}{Версия препринта}.
Роль в плане: карта современной теории.

\bibitem{Satellite}
D.~Dudko, M.~Lyubich. \emph{Local connectivity of the Mandelbrot set
at some satellite parameters of bounded type}.
\href{https://arxiv.org/abs/1808.10425}{arXiv:1808.10425}.
Локальный PDF: \repo{refs/Dudko_1808.10425.pdf}.

\bibitem{Feigenbaum}
D.~Dudko, M.~Lyubich. \emph{MLC at Feigenbaum points}.
\href{https://arxiv.org/abs/2309.02107}{arXiv:2309.02107}.
Локальный PDF: \repo{refs/Dudko_2309.02107.pdf}.

\bibitem{KKL}
J.~Kahn, A.~Kapiamba, M.~Lyubich.
\emph{MLC for parabolically bounded primitive renormalization}.
\href{https://arxiv.org/abs/2606.27272}{arXiv:2606.27272v1}, 2026.
Локальный исходник: \repo{refs/arxiv-2606.27272/waves-with-Misha.tex}.

\bibitem{BDI}
L.~Bartholdi, D.~Dudko. \emph{Algorithmic aspects of branched coverings I.
Van Kampen's Theorem for bisets}.
\href{https://arxiv.org/abs/1512.08539}{arXiv:1512.08539}.
Локальный исходник: \repo{refs/arxiv-1512.08539/bc1-bisets.tex}.

\bibitem{BDIV}
L.~Bartholdi, D.~Dudko. \emph{Algorithmic aspects of branched coverings
IV/V. Expanding maps}.
\href{https://arxiv.org/abs/1610.02434}{arXiv:1610.02434}.
Локальный исходник: \repo{refs/arxiv-1610.02434/bc4-expanding.tex}.

\bibitem{DudkoPage}
D.~Dudko.
\href{https://www.math.stonybrook.edu/~ddudko/Research}{Авторский список работ}.
Источник актуальных названий и версий.

\bibitem{Hatcher}
A.~Hatcher. \emph{Algebraic Topology}. Cambridge University Press, 2002.
\href{https://pi.math.cornell.edu/~hatcher/AT/ATpage.html}{Авторская страница и свободный текст}.

\bibitem{MilnorThird}
J.~Milnor. \emph{Dynamics in One Complex Variable}. Third edition,
Annals of Mathematics Studies 160, Princeton University Press, 2006.
\href{https://www.jstor.org/stable/j.ctt7rnxn}{Издание Princeton University Press}.

\bibitem{MilnorNotes}
J.~Milnor. \emph{Dynamics in One Complex Variable: Introductory Lectures}.
Ранняя версия, revised 1991, \href{https://arxiv.org/abs/math/9201272}{arXiv:math/9201272}.
Местный файл: \repo{refs/9201272v1.pdf}.

\bibitem{NekraIMG}
V.~Nekrashevych. \emph{Iterated monodromy groups}.
В: \emph{Groups St Andrews 2009 in Bath}, vol.~1, LMS Lecture Note Series
387 (2011), pp.~41--93.
\href{https://people.tamu.edu/~nekrash/Preprints/img.pdf}{Предварительная авторская версия, 48 страниц}.

\bibitem{NekraBook}
V.~Nekrashevych. \emph{Self-similar groups}.
Mathematical Surveys and Monographs 117, American Mathematical Society, 2005.
\href{https://people.tamu.edu/~nekrash/publ.html}{Авторский список работ};
\href{https://www.math.stonybrook.edu/~mlyubich/Archive/Stock/Nekrash.pdf}{Ранняя версия монографии}.

\bibitem{KarpickeBlunt}
J.~D.~Karpicke, J.~R.~Blunt.
\emph{Retrieval practice produces more learning than elaborative studying
with concept mapping}. Science 331 (2011), no.~6018, pp.~772--775.
\href{https://learninglab.psych.purdue.edu/downloads/2011/2011_Karpicke_Blunt_Science.pdf}{Текст статьи}.
\end{thebibliography}

\backmatter
\chapter*{Первый рабочий день}
\addcontentsline{toc}{chapter}{Первый рабочий день}
Откройте \repo{novelty_audit.pdf}. Восстановите равенство
\(\psi(H_1)=\{(u,v):\beta(u)=\beta(v)\}\), затем объясните,
почему из него получается рекурсия \(r_m=2r_{m-1}-1\).
Отдельно выпишите, что именно было проверено в Lemma~3.2
и Theorem~3.3 статьи о \(p\)-Basilica-группах.

Во второй сессии начните \repo{templates/novelty_card.md}
для одного следующего вопроса: формулировка, ближайшая теорема,
попытка сведения и предполагаемое добавление.
\end{document}
